\documentclass[11pt,a4paper]{article}
\usepackage[utf8]{inputenc}
\usepackage[T1]{fontenc}
\usepackage{amsmath, amssymb, amsthm}
\usepackage{geometry}
\geometry{margin=2.5cm}
\usepackage{lmodern}
\usepackage[final,expansion=false]{microtype}
\usepackage{booktabs}
\usepackage{array}
\usepackage{enumitem}
\usepackage{xcolor}
\usepackage{titlesec}
\PassOptionsToPackage{hyphens}{url}
\usepackage{hyperref}
\usepackage{xurl}
\hypersetup{
  colorlinks=true,
  linkcolor=blue!50!black,
  urlcolor=blue!50!black,
  citecolor=blue!50!black,
  pdftitle={A dimension-independent collision-depth invariant for right-corner orthoscheme reflection groups},
  pdfauthor={Vico Bonfioli},
}
\usepackage{parskip}
\setlength{\parindent}{0pt}
\setlength{\parskip}{0.55em}

\titleformat{\section}{\normalfont\Large\bfseries\color{blue!40!black}}{\thesection.}{0.6em}{}

\newtheorem{theorem}{Theorem}
\newtheorem{lemma}[theorem]{Lemma}
\newtheorem{corollary}[theorem]{Corollary}
\newtheorem{proposition}[theorem]{Proposition}
\newtheorem{example}[theorem]{Example}
\theoremstyle{remark}
\newtheorem{remark}[theorem]{Remark}
\theoremstyle{plain}

\title{\textbf{A dimension-independent collision-depth invariant\\
for right-corner orthoscheme reflection groups}\\[0.3em]
\large Closed-form growth rate $r_n = 1 + 2\cos(2\pi/(n+3))$ and unconditional Class~C faithfulness in every dimension}
\author{Vico Bonfioli\thanks{Independent researcher, Italy.
\texttt{vicobonfioli@gmail.com}.\,
\url{https://elvec1o.github.io/home/posts/}.}}
\date{\today}

\begin{document}
\sloppy
\maketitle

\begin{abstract}
For every $n \ge 3$ the reflection group $W_n$ generated by the
$n+1$ face-reflections of an $n$-dimensional right-corner orthoscheme
admits a positional combinatorial invariant
$\mathrm{cd}_n(a_1, \ldots, a_n) \in \{\infty, 4, 3\}$, the
\emph{collision-depth}, determined entirely by two integer statistics
$(e, t)$ of the leg sequence: the number of endpoint-adjacent equal
pairs and the number of consecutive equal triples (Theorem~\ref{thm:cd-general}).
The same group has Coxeter--Steinberg growth rate
$r_n = 1 + 2\cos(2\pi/(n+3))$ in closed form (Theorem~\ref{thm:ndim-rate}),
and its affine representation $\rho_n : W_n \to \mathrm{Aff}(\mathbb{R}^n)$
is unconditionally faithful on every leg tuple with pairwise distinct
positive rational coordinates (Theorem~\ref{thm:master-C-uncond}).
The faithfulness proof rests on a uniform Schur-complement determinant
identity $\det Q_n = -\prod_i a_i^2$ (Lemma~\ref{lem:uniform-detQ})
together with rank-$0$ Mordell descent on three Diophantine quartics
whose Jacobians are Cremona curves \texttt{24a1} and \texttt{72a2}.
\end{abstract}

\section{Setting}\label{sec:setting}

Fix $n \ge 2$.  An \emph{$n$-dimensional right-corner orthoscheme}
$O = O(a_1, \ldots, a_n)$ is the simplex in $\mathbb{R}^n$ with
vertex sequence
\[
V_0 = 0, \quad V_k = V_{k-1} + a_k\, e_k \quad (1 \le k \le n),
\]
where $e_1, \ldots, e_n$ is the standard basis and
$(a_1, \ldots, a_n) \in \mathbb{Q}_{>0}^n$.  The $n+1$ bounding
hyperplanes carry reflections $R_0, R_1, \ldots, R_n$; we write
$W_n = W(O) := \langle R_0, \ldots, R_n \rangle \le \mathrm{Isom}(\mathbb{R}^n)$.
The orbit of $O$ under $W_n$ is explored by breadth-first search in
the Cayley graph relative to $\{R_0, \ldots, R_n\}$, and
\[
\mathcal{O}_d(O) = \{ g \cdot O : g \in W_n,\ \text{minimal word-length } d\},
\qquad u_d^n(O) := |\mathcal{O}_d(O)|, \quad u_0 = 1.
\]
All distances are computed in exact rational arithmetic.\footnote{%
A companion paper treats the dimension-$n=2$ case in detail,
including a universality theorem and eight length-$10$ affine
relations specific to planar right triangles, none of which
generalise to $n \ge 3$.  The present paper is self-contained and
does not depend on that companion.}

\paragraph{Independence of legs.}  A symbolic computation in
$\mathbb{Q}(a_1, \ldots, a_n)$ shows that the perpendicular face-pair
structure of $O$ is independent of the leg values: pairs $(i, j)$
with $|i - j| \ge 2$ commute identically, while pairs $(i, i+1)$
have a leg-dependent dihedral angle with $m_{i,i+1} = \infty$ for
generic legs.  The abstract Coxeter group is therefore the
right-angled Coxeter group on the path diagram $P_{n+1}$:
\begin{equation}\label{eq:Wn-presentation}
W_n  \;\cong\;  \bigl\langle\, R_0, \ldots, R_n \,\bigm|\, R_i^2 = 1,\
(R_i R_j)^2 = 1 \text{ for } |i - j| \ge 2 \,\bigr\rangle,
\end{equation}
the same group for every nondegenerate orthoscheme of dimension $n$.
This is the standard right-angled Coxeter system on the commutation
graph that is the complement of $P_{n+1}$.

\paragraph{Steinberg upper bound.}  Steinberg's word-growth formula
(see Humphreys~\cite{humphreys90}, §5.12) gives
$1 / W_n(t^{-1}) = \sum_I (-1)^{|I|} / W_I(t)$, summed over cliques
$I$ of the commutation graph $G_n$ (cliques correspond to subsets
of pairwise-commuting generators, hence to finite parabolics
$W_I \cong (\mathbb{Z}/2)^{|I|}$).  The number of $k$-cliques of
$G_n$ is $\binom{n-k+2}{k}$, so substituting $W_I(t) = (1+t)^{|I|}$:
\begin{equation}\label{eq:steinberg-Wn}
\frac{1}{W_n(t^{-1})}  \;=\;  \sum_{k \ge 0} (-1)^k \binom{n - k + 2}{k}\,
\frac{1}{(1+t)^k}.
\end{equation}
For every faithful or non-faithful affine representation of $W_n$ one
has $u_d^n \le [t^d]\, W_n(t)$.  The aim of \S\S\ref{sec:rate}--\ref{sec:cd}
is (a) to identify $W_n(t)$ in closed form, (b) to certify equality
$u_d^n = [t^d]\, W_n(t)$ unconditionally for distinct positive rational
legs, and (c) to read off the depth at which the inequality first
becomes an equality (the collision-depth).

\section{The closed-form growth rate $r_n$}\label{sec:rate}

\begin{theorem}[Closed-form $n$-dim Coxeter rate]
\label{thm:ndim-rate}
For every $n \ge 2$, the abstract orbit-growth rate of $W_n$ is
\[
r_n  \;=\;  1 \,+\, 2 \cos\!\left(\frac{2\pi}{n+3}\right).
\]
\end{theorem}

\begin{proof}
Start from Steinberg's identity~\eqref{eq:steinberg-Wn}.  Substitute
$y = -1/(1+t)$, so that $(1+t)^{-k} = (-y)^k = (-1)^k y^k$ and the
two $(-1)^k$ factors cancel:
\[
\frac{1}{W_n(t^{-1})}  \;=\;  \sum_{k \ge 0} \binom{n-k+2}{k}\, y^k
                      \;=:\;  f_n(y).
\]
The polynomial $f_n(y)$ is the standard Fibonacci--Pell polynomial:
writing $P_0(y) = 0$, $P_1(y) = 1$, and
$P_m(y) = P_{m-1}(y) + y\, P_{m-2}(y)$ for $m \ge 2$, Pascal's
identity $\binom{n-k+2}{k} = \binom{n-k+1}{k} + \binom{n-k+1}{k-1}$
gives $f_n(y) = P_{n+3}(y)$.  The characteristic equation
$x^2 = x + y$ has roots $\alpha, \beta = (1 \pm \sqrt{1+4y})/2$, and
the Binet formula yields $P_m(y) = (\alpha^m - \beta^m)/(\alpha - \beta)$;
hence $P_m(y) = 0$ iff $(\alpha/\beta)^m = 1$.  Setting
$\alpha/\beta = e^{2\pi i j / m}$ and solving for $y$:
\[
y_j  \;=\;  -\frac{1}{4 \cos^2(\pi j / m)},
\qquad j = 1, 2, \ldots, \lfloor (m-1)/2 \rfloor.
\]
Translating back via $y = -1/(1+t)$ and using
$4 \cos^2\theta - 1 = 1 + 2 \cos(2\theta)$:
\[
t_j  \;=\;  4 \cos^2\!\left(\frac{\pi j}{n+3}\right) - 1
       \;=\;  1 + 2 \cos\!\left(\frac{2\pi j}{n+3}\right),
\qquad 1 \le j \le \lfloor (n+2)/2 \rfloor.
\]
The asymptotic growth rate is the largest root, attained at $j = 1$.\qed
\end{proof}

\begin{corollary}[asymptotic regime]
$r_n \to 3$ as $n \to \infty$, with the explicit residual
\[
3 - r_n  \;=\;  4 \sin^2\!\left(\frac{\pi}{n+3}\right)
         \;\sim\;  \frac{4 \pi^2}{(n + 3)^2}.
\]
\end{corollary}

The closed form gives quadratic-surd values exactly when
$\varphi(n+3) \le 4$ (Euler's totient), namely $n + 3 \in
\{5, 6, 8, 10, 12\}$; see Table~\ref{tab:n-dim-rates}.

\begin{table}[ht!]
\centering
\small
\renewcommand{\arraystretch}{1.15}
\begin{tabular}{c r l}
\toprule
$n$ & $r_n$ & algebraic form \\
\midrule
2  & $1.6180\ldots$ & $\varphi = (1+\sqrt 5)/2$ (golden) \\
3  & $2.0000$       & exactly $2$ \\
4  & $2.2470\ldots$ & $1 + 2\cos(2\pi/7)$, cubic \\
5  & $2.4142\ldots$ & $1 + \sqrt 2$ (silver) \\
6  & $2.5321\ldots$ & $1 + 2\cos(2\pi/9)$, cubic \\
7  & $2.6180\ldots$ & $\varphi^2 = (3+\sqrt 5)/2$ \\
8  & $2.6825\ldots$ & $1 + 2\cos(2\pi/11)$, degree $5$ \\
9  & $2.7321\ldots$ & $1 + \sqrt 3$ \\
10 & $2.7709\ldots$ & $1 + 2\cos(2\pi/13)$, degree $6$ \\
\bottomrule
\end{tabular}
\caption{The closed-form rate $r_n = 1 + 2\cos(2\pi/(n+3))$ for
the right-corner orthoscheme reflection group, verified to $50$-digit
precision for $n$ up to $500$.}
\label{tab:n-dim-rates}
\end{table}

\section{Unconditional Class C faithfulness}\label{sec:classC}

The Steinberg series~\eqref{eq:steinberg-Wn} is an upper bound: only
when the affine representation $\rho_n : W_n \to \mathrm{Aff}(\mathbb{R}^n)$
is faithful and no additional Coxeter relation closes does the BFS
orbit count $u_d^n$ saturate $[t^d]\, W_n(t)$.  Following the
terminology of the 3D analysis, we call a leg tuple
$(a_1, \ldots, a_n)$ \emph{Class C} when its coordinates are
pairwise distinct.  The main result of this section certifies that
the Class C case is unconditional in every dimension.

\subsection{The uniform Schur-complement determinant identity}

The Lorentzian Gram matrix of the orthoscheme in staircase
coordinates --- with simple roots $n_0 = e_1$,
$n_i = a_i e_{i+1} - a_{i+1} e_i$ ($1 \le i \le n-1$), $n_n = e_n$,
and the standard $(n, 1)$ endpoint correction --- is the tridiagonal
matrix $Q_n$ of size $(n+1) \times (n+1)$ given by
\begin{align*}
Q_n[0,0] &= 1 - a_1^2, & Q_n[i,i] &= a_i^2 + a_{i+1}^2\ (1 \le i \le n-1),
& Q_n[n,n] &= 1, \\
Q_n[0,1] &= a_2, & Q_n[i,i+1] &= a_i\, a_{i+2}\ (1 \le i \le n-2),
& Q_n[n-1,n] &= a_{n-1}.
\end{align*}

\begin{lemma}[Uniform Schur-complement determinant identity]
\label{lem:uniform-detQ}
For every $n \ge 3$ and every positive tuple $(a_1, \ldots, a_n)$,
\[
\det Q_n  \;=\;  -\, \prod_{i=1}^{n} a_i^2.
\]
\end{lemma}

\begin{proof}
Let $D_k = \det Q_n[:k{+}1,:k{+}1]$ be the leading principal minor
of order $k+1$.  Tridiagonal expansion along the last row gives
\begin{equation}\tag{$\dagger$}
D_k \;=\; Q_n[k,k]\, D_{k-1} \;-\; Q_n[k-1,k]^2\, D_{k-2}
\qquad (1 \le k \le n),
\end{equation}
with $D_{-1} = 1$.  We claim
\begin{equation}\tag{$\star$}
D_k \;=\; P_k \;:=\;
\Bigl(\prod_{i=1}^k a_i^2\Bigr)\Bigl(1 - \sum_{i=1}^{k+1} a_i^2\Bigr)
\qquad \text{for every } 1 \le k \le n-1.
\end{equation}

\emph{Base case.}  $D_0 = 1 - a_1^2$, and
$D_1 = (1 - a_1^2)(a_1^2 + a_2^2) - a_2^2
     = a_1^2(1 - a_1^2 - a_2^2) = P_1$.

\emph{Interior step (uniform).}  For $2 \le k \le n-1$ the
recurrence ($\dagger$) reads
$D_k = (a_k^2 + a_{k+1}^2)\, D_{k-1} - (a_{k-1} a_{k+1})^2\, D_{k-2}$.
Introduce the abstract symbols
$S_{k-1} := \sum_{i=1}^{k-1} a_i^2$ and
$P_{k-2} := \prod_{i=1}^{k-2} a_i^2$; substituting the closed forms
($\star$) reduces the step to the polynomial identity
\begin{multline*}
P_{k-2}\, a_{k-1}^2\, a_k^2\, (1 - S_{k-1} - a_k^2 - a_{k+1}^2)
\;=\;
(a_k^2 + a_{k+1}^2)\, P_{k-2}\, a_{k-1}^2\, (1 - S_{k-1} - a_k^2) \\
\;-\;
(a_{k-1} a_{k+1})^2\, P_{k-2}\, (1 - S_{k-1})
\end{multline*}
in the five independent symbols $(a_{k-1}, a_k, a_{k+1}, S_{k-1},
P_{k-2})$.  This identity is independent of $k$ and of $n$; SymPy
verifies $\mathrm{expand}(\text{LHS} - \text{RHS}) = 0$ in
$\mathbb{Z}[a_{k-1}, a_k, a_{k+1}, S_{k-1}, P_{k-2}]$.

\emph{Final step (uniform).}  At $k = n$ the recurrence uses
$Q_n[n,n] = 1$ and $Q_n[n-1,n]^2 = a_{n-1}^2$, giving
$D_n = D_{n-1} - a_{n-1}^2\, D_{n-2}$.  With
$S_{n-1} = \sum_{i=1}^{n-1} a_i^2$, $P_{n-2} = \prod_{i=1}^{n-2} a_i^2$,
\[
D_{n-1} - a_{n-1}^2\, D_{n-2}
\;=\; P_{n-2}\, a_{n-1}^2\, \bigl[(1 - S_{n-1} - a_n^2) - (1 - S_{n-1})\bigr]
\;=\; -P_{n-2}\, a_{n-1}^2\, a_n^2
\;=\; -\prod_{i=1}^n a_i^2.
\]
This final identity is again a single polynomial identity in the four
independent symbols $(a_{n-1}, a_n, S_{n-1}, P_{n-2})$, independent of
$n$.  The three identities (base, interior, final) jointly close the
induction for every $n \ge 3$, with no per-$n$ verification required.
\qed
\end{proof}

\subsection{Diophantine emptiness on three Cremona curves}

A non-perpendicular dihedral pair $(i, i+1)$ in the staircase
orthoscheme has $\cos^2 \theta_{i, i+1} \in \mathbb{Q}$; by the
extended Niven theorem (a consequence of Conway--Jones~\cite{conway-jones}:
if $\theta \in \mathbb{Q}\pi$ and $\cos^2\theta \in \mathbb{Q}$, then
$\cos\theta \in \{0, \pm 1/2, \pm\sqrt 2/2, \pm\sqrt 3/2, \pm 1\}$),
any non-trivial closing relation $(R_i R_{i+1})^m = 1$ forces
$\cos^2 \theta_{i, i+1} \in \{0, 1/4, 1/2, 3/4, 1\}$.  The boundary
values $\{0, 1\}$ are immediately incompatible with positive distinct
legs.  The remaining three cut out, in the two legs flanking the
interior pair (with $X = a_i^2$, $Z = a_{i+2}^2$), the Diophantine
quartics
\[
\mathcal{V}_{1/4} :\ X^2 + 14 X Z + Z^2 = m^2, \qquad
\mathcal{V}_{1/2} :\ X^2 + 6 X Z + Z^2 = k^2, \qquad
\mathcal{V}_{3/4} :\ 9 X^2 + 30 X Z + 9 Z^2 = k^2.
\]

\begin{lemma}\label{lem:cremona-empty}
Each of $\mathcal{V}_{1/4}$, $\mathcal{V}_{1/2}$, $\mathcal{V}_{3/4}$
is empty over distinct-coordinate $\mathbb{Q}_{>0}^2$ (i.e., the only
rational solutions have $X Z = 0$ or $X = Z$).
\end{lemma}

\begin{proof}
\emph{$\mathcal{V}_{1/2}$.}  The quartic
$X^2 + 6 X Z + Z^2 = k^2$ is birationally equivalent to the elliptic
curve $E_6 : y^2 = x^3 - x$ (Mordell~\cite{mordell-1969}, §17).
$E_6$ has Mordell--Weil rank $0$ over $\mathbb{Q}$ by Fermat's
classical descent (Cassels~\cite{cassels-1991}, Ch.~II), and its
full rational-point set is the 2-torsion
$\{\infty, (0, 0), (\pm 1, 0)\}$, pulling back to $X Z = 0$ or
$X = Z$.

\emph{$\mathcal{V}_{1/4}$.}  The Jacobian elliptic curve is
\[
E_{1/4} :\ y^2 = x(x - 12)(x - 16),
\]
which is Cremona \texttt{24a1} (conductor $24$, $j = 35\,152/9$,
Mordell--Weil rank $0$ over $\mathbb{Q}$, torsion
$\mathbb{Z}/4 \oplus \mathbb{Z}/2$); see~\cite{cremona-1997}
and the LMFDB record~\cite{lmfdb-24a1}.  Each of its eight rational
points pulls back to $X Z = 0$ or $X = Z$.

\emph{$\mathcal{V}_{3/4}$.}  The quartic $9 X^2 + 30 X Z + 9 Z^2 = k^2$
admits a standard conic-to-Weierstrass substitution and reduces to
\[
E_{3/4} :\ y^2 = x(x - 300)(x - 1200),
\]
Cremona \texttt{72a2} (conductor $72$, same $j$-invariant $35\,152/9$
as $E_{1/4}$ --- they are quadratic twists --- and also rank $0$ with
torsion $\mathbb{Z}/2 \oplus \mathbb{Z}/2$); see~\cite{cremona-1997,
lmfdb-72a2}.  Each torsion point again pulls back to a trivial
quartic solution.\qed
\end{proof}

\subsection{Faithfulness theorem}

\begin{theorem}[Class C faithfulness, all $n \ge 3$, unconditional]
\label{thm:master-C-uncond}
For every $n \ge 3$ and every leg tuple $(a_1, \ldots, a_n) \in
\mathbb{Q}_{>0}^n$ with pairwise distinct coordinates, the affine
representation $\rho_n : W_n \to \mathrm{Aff}(\mathbb{R}^n)$ is
faithful, no Coxeter relation closes beyond the
$\binom{n}{2} - n$ perpendicular pairs, and the BFS layer counts
saturate the abstract Coxeter sphere size:
\[
u_d^n(O) \;=\; [t^d]\, W_n(t)
\qquad \text{for every } d \ge 0.
\]
\end{theorem}

\begin{proof}
By Lemma~\ref{lem:uniform-detQ}, $\det Q_n = -\prod_{i=1}^n a_i^2 \ne 0$
for every positive tuple, so the Lorentzian Coxeter system is
non-degenerate.  Tits' theorem on faithfulness of the geometric
representation (Humphreys~\cite{humphreys90}, Thm.~5.4) then gives
faithfulness of $\rho_n$ on every positive-rational leg tuple.

It remains to exclude additional Coxeter relations.  Each
non-perpendicular pair is of type $(i, i+1)$, $0 \le i \le n-1$;
its cosine satisfies $\cos^2 \theta_{i, i+1} \in \mathbb{Q}$
(boundary pairs $(0,1)$ and $(n-1,n)$ involve two legs, interior
pairs involve three).  By the extended Niven theorem, a closing
Coxeter relation requires
$\cos^2 \theta_{i, i+1} \in \{0, 1/4, 1/2, 3/4, 1\}$.  The boundary
values $\{0, 1\}$ contradict positivity of legs.  At boundary pairs,
the values $1/4, 3/4$ force $a_i = a_{i+1}/\sqrt 3$ or
$a_i = a_{i+1}$, the first irrational and the second excluded by
the distinct-coordinate hypothesis; the value $1/2$ forces
$a_i = a_{i+1}$, again excluded.  At each interior pair, the three
values $1/4, 1/2, 3/4$ define exactly the three Diophantine
quartics $\mathcal{V}_{1/4}, \mathcal{V}_{1/2}, \mathcal{V}_{3/4}$
in the two flanking legs $(a_i, a_{i+2})$ (the dihedral analysis is
uniform in $n$).  By Lemma~\ref{lem:cremona-empty} each of the
three is empty over distinct-coordinate $\mathbb{Q}_{>0}^2$.

Hence outside the symmetry locus $\{a_i = a_{i+1}\}$ for some $i$,
no non-perpendicular dihedral is a rational multiple of $\pi$, no
closing relation appears, and $u_d^n(O) = [t^d]\, W_n(t)$.\qed
\end{proof}

\section{Dimension-independent collision-depth}\label{sec:cd}

For a leg sequence $(a_1, \ldots, a_n)$ define the
\emph{collision-depth} as
\[
\mathrm{cd}_n(a_1, \ldots, a_n)
\;=\; \inf\bigl\{\, d \ge 0 : u_d^n < [t^d]\, W_n(t) \,\bigr\}
\;\in\; \{1, 2, 3, \ldots, \infty\}.
\]
By Theorem~\ref{thm:master-C-uncond}, $\mathrm{cd}_n = \infty$ for
every pairwise-distinct tuple.  The next theorem extends the
description to every leg multiplicity pattern and shows that, in
spite of the multiplicity partition having $p(n)$ possible shapes,
the collision-depth takes only three values, determined by two
local positional statistics.

For $(a_1, \ldots, a_n) \in \mathbb{Q}_{>0}^n$ set
\begin{align*}
e(a_1, \ldots, a_n) &= \#\bigl\{ i \in \{1, n-1\} : a_i = a_{i+1}\bigr\}
\;\in\; \{0, 1, 2\}, \\
t(a_1, \ldots, a_n) &= \#\bigl\{ i : 1 \le i \le n-2,\ a_i = a_{i+1} = a_{i+2}\bigr\}.
\end{align*}
Thus $e$ counts equal pairs at the two \emph{endpoint} (boundary)
positions of the staircase, and $t$ counts \emph{interior triples}
of three consecutive equal legs.  An equal pair sitting strictly in
the interior of the staircase --- such as $a_2 = a_3$ in
$(1, 2, 2, 3)$ at $n = 4$ --- contributes to neither $e$ nor $t$.

\begin{theorem}[Dimension-independent collision-depth, all $n \ge 3$]
\label{thm:cd-general}
For every $n \ge 3$ and every leg sequence
$(a_1, \ldots, a_n) \in \mathbb{Q}_{>0}^n$,
\[
\mathrm{cd}_n(a_1, \ldots, a_n) \;=\;
\begin{cases}
\infty & \text{if } e = t = 0, \\
4 & \text{if } e \ge 1,\ t = 0, \\
3 & \text{if } t \ge 1.
\end{cases}
\]
In particular, $\mathrm{cd}_n \in \{\infty, 4, 3\}$ for every $n$,
and the value depends only on the local positional statistics $(e, t)$.
\end{theorem}

\begin{proof}
The non-perpendicular dihedral pairs of $O = O(a_1, \ldots, a_n)$
are $(R_i, R_{i+1})$, $0 \le i \le n-1$.  Their cosines have the
following uniform form:
\begin{align*}
\cos \theta_{0, 1} &= \frac{a_2}{\sqrt{a_1^2 + a_2^2}}, &
\cos \theta_{n-1, n} &= \frac{a_{n-1}}{\sqrt{a_{n-1}^2 + a_n^2}},
\\
\cos \theta_{i, i+1} &= \frac{a_i\, a_{i+2}}
   {\sqrt{(a_i^2 + a_{i+1}^2)(a_{i+1}^2 + a_{i+2}^2)}}
& &\text{for } 1 \le i \le n-2.
\end{align*}

\emph{Endpoint pair $\Rightarrow \mathrm{cd} = 4$.}  An endpoint
pair $a_1 = a_2$ (or $a_{n-1} = a_n$) makes the boundary cosine
$1/\sqrt 2$, so the dihedral is $\pi/4$ and the Coxeter relation
$(R_0 R_1)^4 = 1$ (resp.\ $(R_{n-1} R_n)^4 = 1$) closes.  The
parabolic $\langle R_0, R_1 \rangle$ then has order $8$ rather than
infinite, and the first collision in the orbit BFS appears at
length $4$, giving $\mathrm{cd}_n = 4$.

\emph{Interior triple $\Rightarrow \mathrm{cd} = 3$.}  An interior
triple $a_i = a_{i+1} = a_{i+2}$ ($1 \le i \le n-2$) makes the
interior cosine $\cos \theta_{i, i+1} = a^2 / (2 a^2) = 1/2$ where
$a := a_i$, so the dihedral is $\pi/3$ and the braid relation
$R_i R_{i+1} R_i = R_{i+1} R_i R_{i+1}$ closes at length $3$.
Since a triple subsumes two pairs, the $\pi/3$ relation appears
before any $\pi/4$ relation that the same triple might also induce
at a boundary position, giving $\mathrm{cd}_n = 3$.

\emph{Interior pair only $\Rightarrow \mathrm{cd} = \infty$.}
Suppose $(e, t) = (0, 0)$, but the sequence is not Class C; then
some interior pair $a_j = a_{j+1}$ exists with $1 \le j \le n-1$,
neither boundary position nor part of a triple.  At each
non-perpendicular pair $(i, i+1)$ the cosine depends on the
flanking legs, which are by hypothesis not all equal.  At the
two interior pairs flanking the equal pair --- $(j-1, j)$ and
$(j+1, j+2)$ (whichever exist) --- the cosine takes the
$a_i a_{i+2} / \sqrt{(a_i^2 + a_{i+1}^2)(a_{i+1}^2 + a_{i+2}^2)}$
form with $(a_i, a_{i+2})$ distinct positive rationals (otherwise
the equal pair would extend to a triple, contradicting $t = 0$).
By the extended Niven theorem and Lemma~\ref{lem:cremona-empty},
no closing relation is possible: each of the candidate values
$1/4, 1/2, 3/4$ for $\cos^2$ defines one of the three Diophantine
quartics, all empty over distinct $\mathbb{Q}_{>0}^2$.  At every
non-flanking pair the flanking-leg analysis is the same.  Hence no
Coxeter relation closes, and by Tits faithfulness on the resulting
non-degenerate Coxeter system (Lemma~\ref{lem:uniform-detQ} plus
Humphreys~\cite{humphreys90}, Thm.~5.4) $u_d^n = [t^d]\, W_n(t)$
for every $d$, i.e., $\mathrm{cd}_n = \infty$.

Combining the three cases proves the tri-classification.\qed
\end{proof}

\begin{remark}[the $R$-rule fails in $n \ge 4$]
\label{rem:R-rule}
A natural but incorrect heuristic is to set $\mathrm{cd}_n$ to a
finite value as soon as some consecutive equal pair appears, i.e.,
as soon as the run-length $R := \max_i \#\{j : a_j = a_i\}$ is at
least $2$.  At $n = 3$ this $R$-rule coincides with the $(e, t)$-rule
because every consecutive equal pair in a $3$-term sequence is
automatically at an endpoint.  At $n \ge 4$ the two rules diverge:
the orthoscheme $(1, 2, 2, 3)$ has an interior equal pair
($R = 2$, but $(e, t) = (0, 0)$) and
Theorem~\ref{thm:cd-general} predicts $\mathrm{cd}_4 = \infty$;
empirically the BFS sequence $1, 5, 14, 33, 75, 169, \ldots$ is
bit-identical to the Class C sequence $1, 5, 14, 33, 75, 169$ at
$(1, 2, 3, 5)$ through every observed depth, falsifying the
$R$-rule.
\end{remark}

\section{Atlas: small dimensions and the $(e, t)$ classes}\label{sec:atlas}

We collect a small atlas of representative leg sequences for
$n = 3, 4, 5, 6$, illustrating each attainable $(e, t)$ class.
The BFS layer counts have been verified bit-perfect against
Theorem~\ref{thm:cd-general} in every case.

\begin{table}[ht!]
\centering
\small
\renewcommand{\arraystretch}{1.2}
\begin{tabular}{c l c c l}
\toprule
$n$ & leg sequence & $(e, t)$ & $\mathrm{cd}_n$ & first six BFS layer counts \\
\midrule
3 & $(2, 3, 5)$       & $(0, 0)$ & $\infty$ & $1, 4, 9, 18, 36, 72$ \\
3 & $(1, 1, 2)$       & $(1, 0)$ & $4$      & $1, 4, 9, 18, 35, 67$ \\
3 & $(1, 1, 1)$       & $(2, 1)$ & $3$      & $1, 4, 9, 17, 28, 42$ \\
\midrule
4 & $(1, 2, 3, 5)$    & $(0, 0)$ & $\infty$ & $1, 5, 14, 33, 75, 169$ \\
4 & $(1, 2, 2, 3)$    & $(0, 0)$ & $\infty$ & $1, 5, 14, 33, 75, 169$ \\
4 & $(1, 1, 2, 3)$    & $(1, 0)$ & $4$      & $1, 5, 14, 33, 73, 161$ \\
4 & $(1, 1, 2, 2)$    & $(2, 0)$ & $4$      & $1, 5, 14, 33, 71, \ldots$ \\
4 & $(1, 1, 1, 2)$    & $(1, 1)$ & $3$      & $1, 5, 14, 31, \ldots$ \\
\midrule
5 & $(1, 2, 3, 5, 7)$    & $(0, 0)$ & $\infty$ & $1, 6, 20, 54, 136, \ldots$ \\
5 & $(1, 2, 2, 3, 5)$    & $(0, 0)$ & $\infty$ & $1, 6, 20, 54, 136, \ldots$ \\
5 & $(1, 1, 2, 3, 5)$    & $(1, 0)$ & $4$      & $1, 6, 20, 54, 133, \ldots$ \\
5 & $(2, 1, 1, 1, 2)$    & $(0, 1)$ & $3$      & $1, 6, 20, 51, \ldots$ \\
\midrule
6 & $(1, 1, 2, 3, 2, 2)$ & $(2, 0)$ & $4$      & $1, 7, 27, 84, 234, \ldots$ \\
\bottomrule
\end{tabular}
\caption{A small atlas of representative orthoschemes in $n = 3, 4, 5, 6$
illustrating each $(e, t)$ class.  The collision-depth value
$\mathrm{cd}_n \in \{\infty, 4, 3\}$ is determined uniformly by
$(e, t)$ in every dimension, by Theorem~\ref{thm:cd-general}.  The
two $n = 4$ sequences with $(e, t) = (0, 0)$ have bit-identical
BFS counts through every observed depth despite differing in run
structure ($(1, 2, 3, 5)$ has $R = 1$ but $(1, 2, 3, 5)$ and $(1, 2, 2, 3)$ have $R = 2$; only $(e, t)$ is the right invariant).}
\label{tab:atlas}
\end{table}

The all-distinct rows of the table saturate the Steinberg upper
bound $[t^d]\, W_n(t)$ by Theorem~\ref{thm:master-C-uncond}; for
$n = 3$ this reproduces the classical sequence
$u_d = 9 \cdot 2^{d-2}$ for $d \ge 2$.

\section{Related work and open questions}\label{sec:related}

\paragraph{Path Coxeter groups and the Cannon--Floyd--Parry framework.}
The closed form $r_n = 1 + 2\cos(2\pi/(n+3))$ is the structural
signature of a linear (path-shaped) Coxeter group whose
Poincar\'e-series denominator reduces to a tridiagonal Chebyshev-type
characteristic polynomial.  Cannon and Wagreich~\cite{cannon-wagreich}
and Floyd and Plotnick~\cite{floyd-plotnick} developed the
growth-rate calculus of such groups; Theorem~\ref{thm:ndim-rate}
fits into their framework with parameter $m = n + 3$ and is
recovered cleanly through the Fibonacci--Pell substitution of
\S\ref{sec:rate}.

\paragraph{Salem-type family on the surface-dynamics side.}
The family of algebraic integers $\{ 1 + 2 \cos(2\pi/m)\}_{m \ge 3}$
to which $r_n$ belongs also arises in birational surface dynamics.
Diller and Favre~\cite{diller-favre} introduced the
dynamical-degree framework for bimeromorphic maps of surfaces;
McMullen~\cite{mcmullen-2007} showed that this family arises as
dynamical degrees of automorphisms of blowups of $\mathbb{P}^2$
associated with the $E_n$ T-shaped Coxeter graph (different from
the path Coxeter graph here); and Bedford--Kim~\cite{bedford-kim}
realised specific members through linear-fractional recurrences.
The rate $r_n$ lies in this family, but the right-corner orthoscheme
reflection group is not known to be realised as a McMullen
involution-product on a rational surface; the coincidence is at
the level of the algebraic-number family.  A birational-geometry
construction realising $r_n$ as a dynamical degree on a
$\mathbb{P}^2$-blowup would be of independent interest.

\paragraph{Open questions.}
(1)~\emph{Non-orthoscheme reflection groups.}  Does an analogue of
Theorem~\ref{thm:cd-general} hold for other families of simplicial
reflection chambers in $\mathbb{R}^n$ --- for example, Coxeter
orthoschemes with bent paths, or simplices arising from
non-path-shaped Coxeter diagrams?  The proof relies on the
right-angled-with-one-$\infty$ structure of the path Coxeter
diagram; how robust is the $(e, t)$-classification under deformation
of the diagram?

(2)~\emph{Classes of intermediate growth.}  In every dimension
$n \ge 3$, partial multiplicity patterns (neither distinct nor
saturated by triples) produce intermediate growth rates strictly
between two consecutive values of $r_n$.  An algebraic
characterisation of these intermediate rates --- in particular,
whether they are algebraic over $\mathbb{Q}$ at all --- is open.

(3)~\emph{Beyond $(e, t)$.}  Theorem~\ref{thm:cd-general} shows that
the collision-depth is determined by the maximum multiplicity of
\emph{positional} structure $(e, t)$.  Is there a finer
positional invariant that captures the entire $u_d^n$ sequence (not
only $\mathrm{cd}_n$) in terms of $(a_1, \ldots, a_n)$?

\section*{Acknowledgments}

Computational assistance from large language model assistants was used
at multiple stages of this work; all mathematical claims have been
verified by the author, who is solely responsible for the content of
this paper and for any errors.

\begin{thebibliography}{99}

\bibitem{humphreys90}
J.~E. Humphreys, \emph{Reflection Groups and Coxeter Groups},
Cambridge Studies in Advanced Mathematics 29, Cambridge University
Press, 1990.

\bibitem{steinberg-1968}
R.~Steinberg, \emph{Endomorphisms of Linear Algebraic Groups},
Memoirs of the American Mathematical Society 80, 1968.

\bibitem{davis-2008}
M.~W. Davis, \emph{The Geometry and Topology of Coxeter Groups},
London Mathematical Society Monographs 32, Princeton University
Press, 2008.

\bibitem{conway-jones}
J.~H. Conway, A.~J. Jones, \emph{Trigonometric Diophantine equations
(On vanishing sums of roots of unity)}, Acta Arith.\ 30 (1976),
229--240.

\bibitem{mordell-1969}
L.~J. Mordell, \emph{Diophantine Equations}, Pure and Applied
Mathematics 30, Academic Press, 1969.

\bibitem{cassels-1991}
J.~W.~S. Cassels, \emph{Lectures on Elliptic Curves}, London
Mathematical Society Student Texts 24, Cambridge University Press,
1991.

\bibitem{cremona-1997}
J.~E. Cremona, \emph{Algorithms for Modular Elliptic Curves},
2nd edition, Cambridge University Press, 1997.

\bibitem{lmfdb-24a1}
LMFDB Collaboration, elliptic curve \texttt{24a1} over $\mathbb{Q}$,
\url{https://www.lmfdb.org/EllipticCurve/Q/24/a/1}.

\bibitem{lmfdb-72a2}
LMFDB Collaboration, elliptic curve \texttt{72a2} over $\mathbb{Q}$,
\url{https://www.lmfdb.org/EllipticCurve/Q/72/a/2}.

\bibitem{cannon-wagreich}
J.~W. Cannon, P.~Wagreich, \emph{Growth functions of surface groups},
Math.\ Ann.\ 293 (1992), no.~2, 239--257.

\bibitem{floyd-plotnick}
W.~J. Floyd, S.~P. Plotnick, \emph{Growth functions on Fuchsian
groups and the Euler characteristic}, Invent.\ Math.\ 88 (1987),
no.~1, 1--29.

\bibitem{parry-1989}
W.~Parry, \emph{Growth series of some hyperbolic Coxeter groups},
Geom.\ Dedicata 32 (1989), 75--96.

\bibitem{diller-favre}
J.~Diller, C.~Favre, \emph{Dynamics of bimeromorphic maps of
surfaces}, Amer.\ J.\ Math.\ 123 (2001), no.~6, 1135--1169.

\bibitem{mcmullen-2007}
C.~T. McMullen, \emph{Dynamics on blowups of the projective plane},
Publ.\ Math.\ Inst.\ Hautes \'Etudes Sci.\ 105 (2007), 49--89.

\bibitem{bedford-kim}
E.~Bedford, K.~Kim, \emph{Dynamics of rational surface automorphisms:
linear fractional recurrences}, J.\ Geom.\ Anal.\ 19 (2009), no.~3,
553--583.

\bibitem{coxeter-1934}
H.~S.~M. Coxeter, \emph{Discrete groups generated by reflections},
Annals of Math.\ 35 (1934), 588--621.

\end{thebibliography}

\end{document}
